\documentclass[10pt,a4paper]{article}
\usepackage[T1]{fontenc}
\usepackage{lmodern,amsmath,booktabs}
\usepackage[margin=22mm]{geometry}
\usepackage[hidelinks]{hyperref}
\setlength{\parindent}{0pt}
\setlength{\parskip}{6pt}
\begin{document}
\begin{center}
{\Large H$\alpha$ contrast investigation: retained results}\\
PR \#16 --- experiments of 27--28 November 2026
\end{center}

\textbf{Disposition.}
The tested effects did not establish a physical explanation of the blue peak
and central trough. The dense outer trial was not adopted. Experimental
drivers, intermediate states, generated figures and verbose verification
receipts are retained in the historical commit, rather than the active tree.

\textbf{Data and comparison.}
The most consistent complete-band target is the first night's twelve
exposures at phase $0.3761496$--$0.3932226$. Bands B/C/R are
$[-40,-16)/[-8,8)/[16,40)$ km s$^{-1}$. The second night lacks complete blue
coverage. Native sampling, masks, common reference and fixed GP covariance
were retained; changed templates reconditioned the GP mean, and the same
global RV and amplitude were profiled for each shape. Operator/master/Doppler
controls found no implementation error in their tested scope. Upstream
extraction and covariance remain unresolved, so these comparisons do not
establish calibrated significance or a unique atmospheric cause.

\textbf{Candidate causes.}
Suppressing the upper source improves the core but leaves a blue deficit.
The tested velocity, longitude and smooth stellar corrections do not repair
both features. Original free model \#1 reduces the target B/C residuals to
about $+184/-165$ ppm, but its full-data $Q$ gain changes from $+21.83$ to
$-0.52$ between pixel selections. Its local success is therefore a structural
clue, not an identified physical mechanism. A fixed-bath density probe looked
promising but did not include radiation, charge and thermal feedback.

\textbf{Coupled dense outer trial.}
The full equations were re-solved with $\rho=q\rho_0$, $v=v_0/q$, using the
current-temperature/composition reference wind and the changed
expansion/compression term. The outer multiplier reached
$D=20$; the required supporting force was prescribed, not derived.
Seventeen roots passed independent CPU checks. None of the seven observed
higher-density cases improved blue-minus-core residuals in either pixel
selection. Primary residuals after the same observation processing are:
\begin{center}
\begin{tabular}{lrrr}
\toprule
Atmosphere & Blue [ppm] & Core [ppm] & Blue$-$core [ppm] \\
\midrule
PR13 reference & $+590.5$ & $-585.5$ & 1175.98 \\
$D=1$ spherical control & $+781.4$ & $-920.2$ & 1701.60 \\
$D=20$ dense trial & $+816.5$ & $-925.2$ & 1741.67 \\
\bottomrule
\end{tabular}
\end{center}
Density itself worsened the primary contrast residual by 40.07 ppm relative
to $D=1$; the entire difference from PR13 must not be attributed to density.
At $1.9375R_p$, $n_2$ increased 332-fold but the source also increased
3.41-fold. The source ratio between $1.5125R_p$ and $1.9375R_p$ was only
1.74, versus about 23 in the free shape. $D=12$ and 15 bracketed the free
model's outer $n_2$ column and still failed the contrast comparison.

\textbf{Limits and archive.}
Target-bracket ray refinement changed blue-minus-core residuals by less than
0.1 ppm. This does not resolve the eight-cell radial structure, physical
support, angular covering or coupled transport adequacy. The result limits
this prescribed density/slow-flow family, not every dense reservoir.
The original implementation, evidence and formal MCEG assessments remain at
\href{https://github.com/HajimeKawahara/exojax-escape/tree/37b277891f5232e9cb8fd1950879952f1adf49d0}{commit \texttt{37b277891f52}}.
The compact historical summary is
\texttt{mceg/investigations/halpha\_contrast.json}.
\end{document}
